% Generated draft: requires mathematical and source review.
\documentclass[11pt]{article}
\usepackage[margin=25mm]{geometry}
\usepackage{fontspec}
\setmainfont{Latin Modern Roman}
\usepackage{amsmath,amssymb,mathtools,unicode-math}
\setmathfont{Latin Modern Math}
\usepackage{enumitem,needspace,xurl,hyperref}
\setlength{\parindent}{0pt}
\setlength{\parskip}{5pt}
\setlength{\emergencystretch}{3em}
\newcommand{\N}{\mathbb{N}}
\newcommand{\Z}{\mathbb{Z}}
\newcommand{\Q}{\mathbb{Q}}
\newcommand{\R}{\mathbb{R}}
\newcommand{\C}{\mathbb{C}}
\newcommand{\F}{\mathbb{F}}
\newcommand{\A}{\mathbb{A}}
\newcommand{\PP}{\mathbb{P}}
\newcommand{\E}{\mathbb{E}}
\newcommand{\eps}{\varepsilon}
\newcommand{\dd}{\,\mathrm{d}}
\begin{document}
\section*{UnsolvedMath ID: 800356}
\textbf{Source problem code:} catalogue entry
\subsection{Definitions and mathematical statement}
Let \(D=\{n\in\Z:n\geq4\}\), and retain the source notation
\[
 f(n)=T_{113}(n,2n-3),\qquad n\in D.
\]
Here \(T_{113}\) refers to the tiling enumeration in the cited source. The supplied extract does not define an \((a,b)\)-benzel, the permitted prototiles, or the meaning of the indices \(113\). Consequently, this notation does not specify a self-contained counting function without the missing source definitions.

Equip \(D\) with the topology induced by the \(2\)-adic topology on \(\Z\): a neighborhood basis at \(n\) consists of the sets \(\{n'\in D:n'\equiv n\pmod{2^r}\}\), for integers \(r\geq0\). For an integer-valued function \(f\), continuity in this topology means
\[
 \forall n\in D\;\forall k\in\Z_{\geq1}\;\exists r\in\Z_{\geq0}\;\forall n'\in D,
 \qquad n'\equiv n\pmod{2^r}
 \Longrightarrow f(n')\equiv f(n)\pmod{2^k}.
\]
\textbf{Question (Propp, 2022, Problem 16).} Does \(f(n)=T_{113}(n,2n-3)\) satisfy this continuity condition? This displays the ordinary topological interpretation of the question; the extract does not specify whether a stronger uniform continuity condition, with \(r\) independent of \(n\), is intended.

\textbf{Supplied observations.} The source reports that the sequence appears to be constant modulo \(2\), constant modulo \(4\), \(2\)-periodic modulo \(8\), and \(8\)-periodic modulo \(16\), with the repeating residue pattern
\[
 2,\ 14,\ 2,\ 14,\ 10,\ 6,\ 10,\ 6.
\]
These are reported observations, not hypotheses or established conclusions in this draft. The extract does not explicitly attach the first displayed residue to an argument. It states that \(n=3\) breaks the pattern; that argument is excluded from \(D\).

\textbf{Additional weighted-trimer context supplied with the record.} At OPAC 2022, David Speyer suggested defining a trimer as a path with three vertices and two edges. According to the source, each stone then corresponds to three trimers, according to which two of its three possible edges are used. Equivalently, one may retain the original trimer definition and give every stone weight \(3\). If \(\mathcal T_{ijk}(a,b)\) denotes the source's class of \(i,j,k\) tilings of the \((a,b)\)-benzel, and \(m(\tau)\) is the number of stones in a tiling \(\tau\), the supplied weighted notation is
\[
 T_{ijk}(a,b;3)=\sum_{\tau\in\mathcal T_{ijk}(a,b)}3^{m(\tau)}.
\]
The extract refers to definitions near the end of Section 1, which are not supplied. It reports that the new numbers have no nice patterns except that, when all prototiles are allowed, stones of weight \(3\) seem to produce the same \(2\)-adic continuity phenomenon as in Problems 11 and 12. Those problems and the precise weighted family under consideration are not supplied. This contextual observation therefore does not determine a separate precise weighted continuity conjecture or change the displayed unweighted question.
\subsection{Short English statement}
Propp asks whether the counts \(T_{113}(n,2n-3)\), for integers \(n\geq4\), vary continuously in the \(2\)-adic topology. Under the ordinary continuity interpretation, sufficiently strong congruence of arguments near each fixed argument must imply any prescribed power-of-two congruence of the counts. The supplied residue patterns are observational. Missing benzel and tile-class definitions prevent a self-contained formulation; the accompanying weight-\(3\) discussion also lacks a precise continuity target.
\subsection{Sources}
\begin{itemize}
\item UnsolvedMath local registry, permanent ID \(800356\), source problem code \texttt{catalogue entry}, title \emph{Trimer Covers in the Triangular Grid --- Problem 16}, in the AMR Open Problem Lists collection. The supplied record attributes the problem to Propp (2022) and labels it \texttt{open}; this is a supplied status label, not an independently verified mathematical or literature assessment.
\item Propp, \emph{Trimer Covers in the Triangular Grid: Twenty Mostly Open Problems} (2022), source item \texttt{Problem 16 (\#8)}: \url{https://arxiv.org/abs/2206.06472}. The supplied provenance records access on 2026-10-16 and extraction by \texttt{arxiv\_latex}. These are supplied provenance details; no independent source visit or verification is claimed. The extract attributes the alternative path definition of trimers to David Speyer at OPAC 2022 and refers to Section 1 and Problems 11 and 12 for additional context.
\item OEIS entry A364417, cited by the supplied statement for the sequence: \url{https://oeis.org/A364417}. Its contents and correspondence with the stated enumeration have not been independently verified.
\end{itemize}
\end{document}
